% Source: https://artofproblemsolving.com/community/c3h2388968_christmas1000th_post_problems
% Stable identity: aops-topic-2388968-q18
Recall that $(\cos x,\sin x)$ is the point reached on the unit circle by turning through the angle $x$ from $(1,0)$.

\textbf{(a)} Turning through $x+2\pi$ is turning through $x$ and then one more full turn, which ends at the same point, so $\sin$ and $\cos$ are unchanged. Turning through $x+\pi$ ends at the diametrically opposite point $(-\cos x,-\sin x)$, so
\[\tan(x+\pi)=\dfrac{-\sin x}{-\cos x}=\dfrac{\sin x}{\cos x}=\tan x.\]

\textbf{(b)} Suppose $T>0$ and $\sin(x+T)=\sin x$ for all $x$. Putting $x=\dfrac\pi2$ gives $\sin\left(\dfrac\pi2+T\right)=1$. The only point of the unit circle with second coordinate $1$ is $(0,1)$, so $\dfrac\pi2+T=\dfrac\pi2+2\pi k$ for an integer $k$; as $T>0$, $k\ge1$ and $T\ge2\pi$. For $\cos$, put $x=0$: $\cos T=1$ forces the point $(1,0)$, so again $T=2\pi k\ge2\pi$. For $\tan$: if $\tan T$ is defined, putting $x=0$ gives $\tan T=0$, so $\sin T=0$ and $T$ is a multiple of $\pi$; if $\tan T$ is not defined, then $T=\dfrac\pi2+k\pi$ and $x=\dfrac\pi4$ gives $\tan\left(\dfrac{3\pi}{4}+k\pi\right)=-1\ne1=\tan\dfrac\pi4$, which is impossible. So $T$ is a positive multiple of $\pi$ and $T\ge\pi$. Together with (a), the smallest periods are $2\pi$, $2\pi$ and $\pi$. $\blacksquare$

\begin{asy}
usepackage("amsmath");
size(7cm);
defaultpen(fontsize(10pt));
import graph;
real xs=0.5;
draw((-0.3,0)--(xs*4*pi+0.4,0),Arrow(6)); draw((0,-1.5)--(0,1.6),Arrow(6));
path s=graph(new pair(real t){return (xs*t,sin(t));},0,4*pi);
draw(s,blue+linewidth(0.9));
for(int k=1;k<=4;++k){draw((xs*k*pi,-0.06)--(xs*k*pi,0.06));}
label("$\pi$",(xs*pi,0),SW,fontsize(8pt)); label("$2\pi$",(xs*2*pi,0),SW,fontsize(8pt)); label("$3\pi$",(xs*3*pi,0),SW,fontsize(8pt)); label("$4\pi$",(xs*4*pi,0),SW,fontsize(8pt));
label("$y=\sin x$",(xs*3.5*pi,-1),S,blue+fontsize(8pt));
draw((xs*pi/2,1.2)--(xs*5*pi/2,1.2),red,Arrows(4)); label("$2\pi$",(xs*1.5*pi,1.2),N,red+fontsize(8pt));
real ys=0.25;
pair o=(0,-4.2);
draw(o+(-0.3,0)--o+(xs*2*pi+0.4,0),Arrow(6)); draw(o+(0,-1.5)--o+(0,1.6),Arrow(6));
for(int k=0;k<=2;++k){
  real a=k*pi-pi/2+0.25, c=k*pi+pi/2-0.25;
  if(k==0) a=0;
  if(k==2) c=2*pi;
  draw(graph(new pair(real t){return o+(xs*t,ys*tan(t));},a,c),heavygreen+linewidth(0.9));
}
draw(o+(xs*pi/2,-1.5)--o+(xs*pi/2,1.5),dashed+gray); draw(o+(xs*3*pi/2,-1.5)--o+(xs*3*pi/2,1.5),dashed+gray);
label("$\pi$",o+(xs*pi,0),SE,fontsize(8pt)); label("$2\pi$",o+(xs*2*pi,0),SE,fontsize(8pt));
label("$y=\tan x$",o+(xs*2*pi,1.2),E,heavygreen+fontsize(8pt));
draw(o+(xs*0,-1.8)--o+(xs*pi,-1.8),red,Arrows(4)); label("$\pi$",o+(xs*pi/2,-1.8),S,red+fontsize(8pt));
\end{asy}

The graph of $\sin x$ repeats every $2\pi$; the graph of $\tan x$ (vertical scale compressed) repeats every $\pi$, between its vertical asymptotes.
